% Part: computability
% Chapter: recursive-functions
% Section: introduction

\documentclass[../../../include/open-logic-section]{subfiles}

\begin{document}

\olfileid[hi]{cmp}{rec}{int}
\olsection{परिचय}

संगणनीयता का गणितीय सिद्धांत विकसित करने के लिए सबसे पहले संगणनीयता का
एक \emph{प्रतिरूप} विकसित करना पड़ता है। आज हम संगणनीयता को उस प्रकार की
क्रिया समझते हैं जो संगणक करते हैं, और संगणक चिह्नों के साथ काम करते हैं।
किंतु संगणनीयता के सिद्धांतों के आरंभिक विकास में संगणन का आदर्श उदाहरण
\emph{संख्यात्मक} संगणन था। गणितज्ञों की रुचि सदा संख्या-सैद्धांतिक फलनों
में रही है, अर्थात् ऐसे फलनों $f\colon \Nat^n \to
\Nat$ में जिनकी गणना की जा सकती है। इसलिए यह आश्चर्यजनक नहीं कि
संगणनीयता-सिद्धांत के आरंभ में ऐसे ही फलनों का अध्ययन हुआ। परिचित
संगणनीय संख्यात्मक फलनों के उदाहरणों---जैसे प्राकृतिक संख्याओं का जोड़,
गुणा और घातांक---में एक रोचक विशेषता समान है: उन्हें \emph{पुनरावर्ती रूप
से} परिभाषित किया जा सकता है। अतः पुनरावर्ती परिभाषाओं के आधार पर
\emph{संगणनीय फलन} की सामान्य परिभाषा देने का प्रयास स्वाभाविक है।
संख्या-सैद्धांतिक फलनों को पुनरावर्ती रूप से परिभाषित करने के अनेक संभव
ढंगों में से परिभाषा का एक विशेष रूप से सरल प्रतिरूप यहाँ केंद्रीय बनता
है: तथाकथित \emph{आदिम पुनरावर्तन}।

संगणनीय फलनों के अतिरिक्त हमारी रुचि संगणनीय समुच्चयों और संबंधों में भी
हो सकती है। कोई समुच्चय संगणनीय है यदि हम यह उत्तर संगणित कर सकें कि दी
हुई संख्या उस समुच्चय की !!a{element} है या नहीं; और कोई संबंध संगणनीय
तभी और केवल तभी है जब हम यह संगणित कर सकें कि टुपल
$\tuple{n_1, \dots, n_k}$ उस संबंध की !!a{element} है या नहीं। किसी
समुच्चय अथवा संबंध के \emph{अभिलाक्षणिक फलन} पर विचार करके संगणनीय
समुच्चयों और संबंधों की चर्चा को संगणनीय फलनों की चर्चा में समाहित किया
जा सकता है। इस प्रकार हम आदिम पुनरावर्ती संबंध भी परिभाषित कर सकते हैं;
मसलन, संबंध ``$n$, $m$ को पूर्णतः विभाजित करता है'' एक आदिम पुनरावर्ती
संबंध है।

फिर भी आदिम पुनरावर्ती फलन---वे जिन्हें केवल आदिम पुनरावर्तन का उपयोग
करके परिभाषित किया जा सकता है---सभी संगणनीय संख्या-सैद्धांतिक फलन नहीं
हैं। आदिम पुनरावर्तन के अनेक सामान्यीकरणों पर विचार किया गया है, किंतु
फलनों को संगणित करने का सबसे शक्तिशाली और व्यापक रूप से स्वीकृत अतिरिक्त
ढंग अपरिबद्ध खोज है। इससे \emph{आंशिक पुनरावर्ती फलनों} की परिभाषा और
उससे संबंधित \emph{सामान्य पुनरावर्ती फलनों} की परिभाषा मिलती है। सामान्य
पुनरावर्ती फलन संगणनीय और पूर्ण होते हैं, और यह परिभाषा ठीक उन आंशिक
पुनरावर्ती फलनों का अभिलक्षणन करती है जो पूर्ण होते हैं। पुनरावर्ती फलन
संगणन के प्रत्येक अन्य प्रतिरूप (ट्यूरिंग मशीन, लैम्ब्डा कलन इत्यादि) का
अनुकरण कर सकते हैं और इसलिए संगणन के अनेक स्वीकृत प्रतिरूपों में से एक
का निरूपण करते हैं।


\end{document}
